\paragraph{Stylized facts and parametric volatility.}
The canonical set of stylized facts \cite{cont2001empirical}, namely heavy tails \cite{mandelbrot1963variation}, absence of linear autocorrelation \cite{fama1970efficient}, and slow decay of the absolute-return ACF \cite{schwert1989does, ding1993long, andersen1997heterogeneous}, must be reproduced jointly for a generator to be useful downstream.
Autoregressive conditional heteroscedasticity (ARCH) and its generalised variant (GARCH) \cite{engle1982autoregressive, bollerslev1986generalized} encode volatility persistence through a single decay parameter but impose a unimodal innovation distribution and, as we confirm, fail distributional tests on equity returns.
Jump-diffusion models \cite{merton1976option, kou2002jump} generate heavy tails but lack a clustering mechanism: after a jump the conditional variance returns immediately to baseline; \citet{maheu2004news} couple jumps with a regime-switching volatility component to address this, an alternative route to regime-dependent tails that our CHMM-t scaffold addresses through heavy-tailed emissions rather than an explicit jump process.
Maximum-likelihood fitting of heavy-tailed innovations in a latent-volatility framework has been studied for stochastic-volatility-in-mean models by \citet{abantovalle2017svm}, who compare Normal, Student-t, generalized error distribution (GED), and Slash errors; our CHMM-t and CHMM-L are the HMM analog of that emission-family choice.

\paragraph{Regime-switching and the Ryd\'{e}n et al.\ limitation.}
\citet{hamilton1989new} introduced Markov-switching models to economics, and \citet{schaller1997regime} provided an early demonstration that monthly equity returns from the Center for Research in Security Prices (CRSP) are well-described by a two-state Markov-switching specification with strongly different mean and variance between regimes; \citet{angbekaert2002regime} subsequently showed that short-rate dynamics in several developed markets also admit a regime-switching representation, establishing regime-switching as a cross-asset-class phenomenon rather than one confined to macro series or equities.
The seminal evaluation of HMMs against the Cont stylized facts is due to \citet{ryden1998stylized}, who fitted Gaussian-emission HMMs with $K = 2$--$3$ states to daily equity-index return series and showed that the geometric sojourn-time distribution of a standard Markov chain produces exponential ACF decay too fast to match the observed slow decay.
\citet{bulla2006stylized} restored fidelity with hidden semi-Markov models (HSMMs) that replace the geometric sojourn distribution with flexible alternatives.
\citet{alswaidan2026hybrid} instead kept the discrete HMM framework and introduced a Poisson jump-duration process to force dwell time in tail states.

Within the same family, \citet{nguyen2018hidden} applies a four-state continuous HMM to monthly S\&P~500 prices for OoS trading signal generation, selecting the state count through the AIC, BIC, HQC, and CAIC information criteria.
We adopt the same four-criterion selection (Section~\ref{sec:model_selection}), in addition to our multi-metric distributional score, to place the $K$-sweep recommendation on standard statistical footing.

Consistent with this broader regime-switching literature, we find that at $K \geq 3$ with quantile-based initialization the Baum-Welch algorithm learns heterogeneous self-transition probabilities across regimes, and the resulting mixture of $K$ geometric sojourn-time distributions approximates slow ACF decay.
This moderate-$K$ route complements the HSMM \citep{bulla2006stylized} and jump-augmented \citep{alswaidan2026hybrid} corrections to the Ryd\'{e}n et al.\ result rather than superseding them; the integrative contribution of this paper is the unified log-space EM scaffold shared by the three emission variants, the twelve-generator evaluation panel, and the regime-conditional VaR test (Section~\ref{sec:var_backtest}), not a single-route resolution of the Ryd\'{e}n limitation.

\paragraph{Discrete HMM with jumps.}
The immediate predecessor of this paper, \citet{alswaidan2026hybrid}, discretizes returns into Laplace quantile bins, estimates $\mathbf{T}$ by frequency counting, fits bin-conditional Student-t emissions, and adds a Poisson jump-duration process with two grid-calibrated hyperparameters ($\epsilon, \lambda$).
The continuous model presented here offers a simpler alternative on each of three axes: no binning, a single EM step in which Baum-Welch jointly learns transitions and emissions (instead of a separate frequency-counted $\mathbf{T}$ plus within-bin emission fit), and no jump mechanism at moderate $K$; this is a modelling-cost comparison rather than a claim that the discrete-plus-jumps route is inferior for the tasks it was designed for.

\paragraph{Cross-asset dependence.}
For multi-asset synthesis, the Single Index Model \cite{sharpe1963simplified} propagates a market factor through a rank-one linear decomposition but distorts each non-market marginal; Sklar's theorem \cite{sklar1959fonctions, nelsen2006introduction} supports a cleaner decomposition in which each asset's fitted marginal is kept and only the copula is estimated.
Gaussian and Student-t copulas \cite{demarta2005tcopula, mcneil2015quantitative} are the standard choices, and rank-reordering simulation \cite{iman1982distribution} injects the copula dependence onto pre-simulated marginal paths without distorting any marginal.
We confirm, on continuous CHMM marginals, the copula-ahead-of-SIM ordering reported by \citet{alswaidan2026hybrid} for discrete HMM marginals, which suggests that the ordering is intrinsic to the rank-reordering construction rather than specific to the marginal model.

\paragraph{Deep generative models and quality assessment.}
Generative adversarial network (GAN) based generators reproduce the visual appearance of return series but typically fail the clustering test unless architectures encode temporal dependence explicitly \cite{takahashi2019modeling, kwon2024can}, the TimeGAN family of \citet{yoon2019timegan} being the standard reference in this direction; neural baselines in \citet{alswaidan2026hybrid} close most of the ACF gap but suffer variance collapse, indicating that the distribution--clustering tension remains in tension for neural approaches.
Following the GRU baseline used in \citet{alswaidan2026hybrid}, we add a single-layer GRU + Gaussian-head auto-regressive recurrent generator (Section~\ref{sec:benchmarks}) to Table~\ref{tab:model_comparison} as a deep-generative head-to-head row, trained by negative log-likelihood on the IS series and rolled forward by sampling from the predicted Gaussian and feeding the sample back into the recurrence.
For evaluation, we adopt the category framework of \citet{stenger2024thinking} and the financial-synthesis emphasis of \citet{assefa2020generating}: no single metric is decisive, so we report a seven-metric panel spanning distributional, tail, and temporal fidelity.
